\documentclass[9pt,a4paper]{article}
\usepackage[a4paper,left=14mm,right=14mm,top=12mm,bottom=12mm,headheight=14pt,headsep=5mm,footskip=7mm]{geometry}
\usepackage{fontspec}
\usepackage{microtype}
\usepackage{array,tabularx,booktabs,longtable}
\usepackage{fancyhdr}
\usepackage{xcolor}
\usepackage{enumitem}
\usepackage{lastpage}
\usepackage{hyperref}
\hypersetup{hidelinks}

\setmainfont{Noto Serif Devanagari}[Script=Devanagari,Language=Hindi]
\newfontfamily\latinfont{Noto Serif}[Ligatures=TeX]
\newfontfamily\monofont{Noto Sans Mono}
\definecolor{ink}{HTML}{202A33}
\definecolor{rule}{HTML}{65727E}
\definecolor{paperblue}{HTML}{294B63}
\definecolor{placeholder}{HTML}{5B6670}
\definecolor{placeholderbg}{HTML}{EEF1F3}
\setlength{\parindent}{0pt}
\setlength{\parskip}{2pt}
\renewcommand{\arraystretch}{1.08}
\setlist[enumerate]{leftmargin=16pt,itemsep=2pt,topsep=2pt}

% All dynamic fields use stable English keys inside [[...]] tokens.
% An AI/data-population process should replace the exact tokens before issuance.
\newcommand{\PH}[1]{\textcolor{placeholder}{\fboxsep=1pt\colorbox{placeholderbg}{\monofont\scriptsize\detokenize{[[#1]]}}}}
\newcommand{\SEC}[1]{\vspace{4pt}{\color{paperblue}\large\bfseries #1}\par\vspace{2pt}\hrule height 0.55pt\vspace{3pt}}
\newcommand{\NOTE}[1]{\textcolor{rule}{\footnotesize\textit{#1}}}

\pagestyle{fancy}
\fancyhf{}
\fancyhead[L]{\small\textbf{प्रोजेक्ट अनंत} \;|\; ऑपरेशन अभेद्य-चक्र}
\fancyhead[R]{\small\textit{गोपनीय — जांच अभिलेख}}
\fancyfoot[L]{\small वित्तीय साइबर अपराध जांच रिपोर्ट}
\fancyfoot[R]{\small पृष्ठ \thepage\ / \pageref{LastPage}}
\renewcommand{\headrulewidth}{0.5pt}
\renewcommand{\footrulewidth}{0.4pt}

\begin{document}

\begin{center}
{\fontsize{13}{15}\selectfont\bfseries वित्तीय साइबर अपराध}\par
{\fontsize{13}{15}\selectfont\bfseries लेनदेन श्रृंखला एवं सिंडिकेट विश्लेषण रिपोर्ट}\par
{\small प्रोजेक्ट अनंत — विश्लेषण रिपोर्ट पहचानकर्ता: \PH{REPORT_ID}}\par
\vspace{3pt}\rule{0.92\linewidth}{0.7pt}
\end{center}

\SEC{1. प्रकरण / घटना सारांश}
\begin{tabularx}{\linewidth}{@{}>{\bfseries}p{29mm} X >{\bfseries}p{29mm} X@{}}
प्राथमिकी क्रमांक & \PH{FIR_NUMBER} & प्रकरण क्रमांक & \PH{CASE_NUMBER}\\
अधिकारी का नाम & \PH{OFFICER_NAME} & थाना & \PH{POLICE_STATION}\\
जिला & \PH{DISTRICT} & दिनांक & \PH{REPORT_DATE}\\
रिपोर्ट पहचानकर्ता & \PH{REPORT_ID} & जांच अवधि & \PH{INVESTIGATION_PERIOD}
\end{tabularx}

\vspace{3pt}
\begin{tabularx}{\linewidth}{@{}>{\bfseries}p{30mm} X >{\bfseries}p{30mm} X@{}}
पीड़ित खाता & \PH{VICTIM_ACCOUNT_MASKED} & बैंक & \PH{VICTIM_BANK}\\
धोखाधड़ी राशि & \PH{FRAUD_AMOUNT} & प्रारंभिक UTR & \PH{INITIAL_UTR}\\
प्रारंभिक लेनदेन & \PH{INITIAL_TRANSACTION_DATETIME} & जांच स्रोत & \PH{INVESTIGATION_SOURCE}
\end{tabularx}

\vspace{3pt}
\textbf{घटना सारांश।} \PH{INCIDENT_SUMMARY}

\vspace{3pt}
\textbf{प्रमुख निष्कर्ष}
\scriptsize
\begin{tabularx}{\linewidth}{@{}l r@{\hspace{14mm}}l r@{}}
विश्लेषित लेनदेन & \PH{TRANSACTIONS_ANALYSED} & विश्लेषित खाते & \PH{ACCOUNTS_ANALYSED}\\
संदिग्ध म्यूल खाते & \PH{SUSPECTED_MULES} & सिंडिकेट & \PH{SYNDICATES_IDENTIFIED}\\
अनुसरित कुल राशि & \PH{TOTAL_TRACED_AMOUNT} & अधिकतम हॉप & \PH{MAXIMUM_HOP}
\end{tabularx}
\normalsize

\SEC{धनराशि प्रवाह}
\NOTE{जांचाधीन लेनदेन श्रृंखला से संबंधित पुनर्निर्मित प्रवाह। तालिका ही मूल अभिलेख है; कोई ग्राफिकल नेटवर्क प्रस्तुति उपयोग नहीं की गई है।}

\scriptsize
\begin{tabularx}{\linewidth}{@{}c l l l r c l@{}}
\toprule
\textbf{हॉप} & \textbf{समय} & \textbf{प्रेषक} & \textbf{प्राप्तकर्ता} & \textbf{राशि} & \textbf{स्तर} & \textbf{UTR / संदर्भ}\\
\midrule
1 & \PH{TX1_TIME} & \PH{TX1_FROM} & \PH{TX1_TO} & \PH{TX1_AMOUNT} & \PH{TX1_LAYER} & \PH{TX1_UTR}\\
2 & \PH{TX2_TIME} & \PH{TX2_FROM} & \PH{TX2_TO} & \PH{TX2_AMOUNT} & \PH{TX2_LAYER} & \PH{TX2_UTR}\\
3 & \PH{TX3_TIME} & \PH{TX3_FROM} & \PH{TX3_TO} & \PH{TX3_AMOUNT} & \PH{TX3_LAYER} & \PH{TX3_UTR}\\
4 & \PH{TX4_TIME} & \PH{TX4_FROM} & \PH{TX4_TO} & \PH{TX4_AMOUNT} & \PH{TX4_LAYER} & \PH{TX4_UTR}\\
5 & \PH{TX5_TIME} & \PH{TX5_FROM} & \PH{TX5_TO} & \PH{TX5_AMOUNT} & \PH{TX5_LAYER} & \PH{TX5_UTR}\\
\bottomrule
\end{tabularx}
\normalsize

\NOTE{केवल जांचाधीन प्रवाह के लिए आवश्यक लेनदेन पंक्तियां रखें। असंबंधित लेनदेन शामिल न करें।}

\newpage
\SEC{2. संदिग्ध खाते का विश्लेषण}
\scriptsize
\begin{tabularx}{\linewidth}{@{}p{17mm} p{25mm} c p{22mm} p{21mm} p{21mm} X@{}}
\toprule
\textbf{खाता} & \textbf{बैंक} & \textbf{स्तर} & \textbf{जोखिम सूचकांक} & \textbf{आवक} & \textbf{जावक} & \textbf{आवक / जावक समकक्ष}\\
\midrule
\PH{A1_ACCOUNT} & \PH{A1_BANK} & \PH{A1_LAYER} & \PH{A1_RISK} & \PH{A1_INFLOW} & \PH{A1_OUTFLOW} & \PH{A1_IN_CTRP} / \PH{A1_OUT_CTRP}\\
\PH{A2_ACCOUNT} & \PH{A2_BANK} & \PH{A2_LAYER} & \PH{A2_RISK} & \PH{A2_INFLOW} & \PH{A2_OUTFLOW} & \PH{A2_IN_CTRP} / \PH{A2_OUT_CTRP}\\
\PH{A3_ACCOUNT} & \PH{A3_BANK} & \PH{A3_LAYER} & \PH{A3_RISK} & \PH{A3_INFLOW} & \PH{A3_OUTFLOW} & \PH{A3_IN_CTRP} / \PH{A3_OUT_CTRP}\\
\bottomrule
\end{tabularx}
\normalsize

\vspace{3pt}\textbf{विश्लेषणात्मक संकेतक}\par
\scriptsize
\begin{tabularx}{\linewidth}{@{}X r@{}}
\toprule
\textbf{संकेतक} & \textbf{मान}\\
\midrule
टर्मिनल कैश-आउट & \PH{INDICATOR_TERMINAL_CASHOUT}\\
साइबर ऑटोमेशन & \PH{INDICATOR_CYBER_AUTOMATION}\\
टर्नओवर संरक्षण & \PH{INDICATOR_TURNOVER_CONSERVATION}\\
समयगत वेग & \PH{INDICATOR_TEMPORAL_VELOCITY}\\
निष्क्रियता-विस्फोट & \PH{INDICATOR_DORMANCY_BURST}\\
समकक्ष असममिति & \PH{INDICATOR_COUNTERPARTY_ASYMMETRY}\\
संरचनात्मक फैन पैटर्न & \PH{INDICATOR_STRUCTURAL_FAN_PATTERN}\\
\bottomrule
\end{tabularx}
\normalsize

\vspace{3pt}\textbf{अवलोकित व्याख्या।} \PH{ACCOUNT_INTERPRETATION}

\SEC{सिंडिकेट विश्लेषण}
\begin{tabularx}{\linewidth}{@{}>{\bfseries}p{25mm} X >{\bfseries}p{25mm} X@{}}
सिंडिकेट पहचानकर्ता & \PH{SYNDICATE_ID} & प्रारूप & \PH{SYNDICATE_ARCHETYPE}\\
सदस्य & \PH{SYNDICATE_MEMBERS} & कुल मात्रा & \PH{SYNDICATE_TOTAL_VOLUME}\\
समय अवधि & \PH{SYNDICATE_TIME_WINDOW} & संबंध & \PH{SYNDICATE_RELATIONSHIP}
\end{tabularx}

\vspace{3pt}
\scriptsize
\begin{tabularx}{\linewidth}{@{}l l c r@{}}
\toprule
\textbf{खाता} & \textbf{भूमिका} & \textbf{स्तर} & \textbf{जोखिम सूचकांक}\\
\midrule
\PH{MEMBER_1_ACCOUNT} & \PH{MEMBER_1_ROLE} & \PH{MEMBER_1_LAYER} & \PH{MEMBER_1_RISK_INDEX}\\
\PH{MEMBER_2_ACCOUNT} & \PH{MEMBER_2_ROLE} & \PH{MEMBER_2_LAYER} & \PH{MEMBER_2_RISK_INDEX}\\
\PH{MEMBER_3_ACCOUNT} & \PH{MEMBER_3_ROLE} & \PH{MEMBER_3_LAYER} & \PH{MEMBER_3_RISK_INDEX}\\
\PH{MEMBER_4_ACCOUNT} & \PH{MEMBER_4_ROLE} & \PH{MEMBER_4_LAYER} & \PH{MEMBER_4_RISK_INDEX}\\
\PH{MEMBER_5_ACCOUNT} & \PH{MEMBER_5_ROLE} & \PH{MEMBER_5_LAYER} & \PH{MEMBER_5_RISK_INDEX}\\
\bottomrule
\end{tabularx}
\normalsize

\NOTE{विश्लेषणात्मक वर्गीकरण केवल देखे गए लेनदेन व्यवहार को दर्शाता है और किसी कानूनी निष्कर्ष का निर्धारण नहीं करता।}

\SEC{3. साक्ष्य सारांश}
\scriptsize
\begin{tabularx}{\linewidth}{@{}l l X X@{}}
\toprule
\textbf{साक्ष्य पहचानकर्ता} & \textbf{प्रकार} & \textbf{संदर्भ} & \textbf{संबंधित निष्कर्ष}\\
\midrule
\PH{EVIDENCE_1_ID} & \PH{EVIDENCE_1_TYPE} & \PH{EVIDENCE_1_REFERENCE} & \PH{EVIDENCE_1_FINDING}\\
\PH{EVIDENCE_2_ID} & \PH{EVIDENCE_2_TYPE} & \PH{EVIDENCE_2_REFERENCE} & \PH{EVIDENCE_2_FINDING}\\
\PH{EVIDENCE_3_ID} & \PH{EVIDENCE_3_TYPE} & \PH{EVIDENCE_3_REFERENCE} & \PH{EVIDENCE_3_FINDING}\\
\PH{EVIDENCE_4_ID} & \PH{EVIDENCE_4_TYPE} & \PH{EVIDENCE_4_REFERENCE} & \PH{EVIDENCE_4_FINDING}\\
\PH{EVIDENCE_5_ID} & \PH{EVIDENCE_5_TYPE} & \PH{EVIDENCE_5_REFERENCE} & \PH{EVIDENCE_5_FINDING}\\
\bottomrule
\end{tabularx}
\normalsize

\vspace{3pt}\textbf{जांच संबंधी निष्कर्ष}
\begin{enumerate}
\item \PH{FINDING_1}
\item \PH{FINDING_2}
\item \PH{FINDING_3}
\item \PH{FINDING_4}
\item \PH{FINDING_5}
\end{enumerate}

\textbf{निष्कर्ष।} \PH{CONCLUSION}

\vspace{4pt}\textbf{अधिकारी प्रमाणन}

मैं प्रमाणित करता/करती हूँ कि यह रिपोर्ट संदर्भित अभिलेखों से पहचाने गए विश्लेषणात्मक अवलोकनों एवं लेनदेन संबंधों को दर्ज करती है तथा मूल साक्ष्यों के विरुद्ध सत्यापन के अधीन है।

\vspace{7pt}
\begin{tabularx}{\linewidth}{@{}X X@{}}
\rule{0.82\linewidth}{0.4pt} & \rule{0.82\linewidth}{0.4pt}\\[-1pt]
\PH{CERTIFYING_OFFICER_NAME} & हस्ताक्षर\\[7pt]
\rule{0.82\linewidth}{0.4pt} & \rule{0.82\linewidth}{0.4pt}\\[-1pt]
\PH{CERTIFYING_OFFICER_RANK} & \PH{CERTIFICATION_DATE}\\[7pt]
\rule{0.82\linewidth}{0.4pt} & \\
\PH{CERTIFYING_POLICE_STATION} & \\
\end{tabularx}

\vfill
\begin{center}
\footnotesize\textbf{दस्तावेज स्थिति: जांच संबंधी विश्लेषणात्मक अभिलेख।}\quad\textit{सभी [[...]] टोकन डेटा प्लेसहोल्डर हैं और जारी करने से पहले बदले जाने चाहिए।}
\end{center}

\end{document}
